% ============================================================
% DOCUMENTCLASS
% ============================================================
\documentclass[12pt,oneside]{book}

% ============================================================
% METADATOS
% ============================================================
\newcommand{\universidad}{Universidad de Buenos Aires (UBA)}
\newcommand{\materia}{Derecho Notarial y Contratos}
\newcommand{\titulo}{El contrato de fideicomiso: estructura, sujetos y aplicaciones prácticas}
\newcommand{\autor}{Isaac Edgar Camacho Ocampo}
\newcommand{\anio}{2026}

% ============================================================
% PAQUETES
% ============================================================
\usepackage[utf8]{inputenc}
\usepackage[T1]{fontenc}
\usepackage[spanish,es-nodecimaldot]{babel}

\usepackage[
    top=2.5cm,
    bottom=2.5cm,
    left=3cm,
    right=2.5cm,
    headheight=15pt
]{geometry}

\usepackage{amsmath}
\usepackage{amssymb}
\usepackage{mathtools}

\usepackage{graphicx}
\usepackage{float}
\usepackage{caption}
\usepackage{subcaption}

\usepackage{booktabs}
\usepackage{array}
\usepackage{longtable}
\usepackage{tabularx}

\usepackage{enumitem}

\usepackage{xcolor}
\usepackage[most]{tcolorbox}

\usepackage{fancyhdr}
\usepackage{ragged2e}

\usepackage[
    colorlinks=true,
    linkcolor=blue!50!black,
    citecolor=green!40!black,
    urlcolor=blue!60!black,
    pdftitle={\titulo},
    pdfauthor={\autor},
    pdfsubject={Apuntes universitarios},
    bookmarks=true
]{hyperref}

% ============================================================
% CONFIGURACIÓN
% ============================================================
\graphicspath{
    {./img/}
    {./graficos/}
    {./figuras/}
}

\setcounter{tocdepth}{2}

\setlength{\parindent}{0pt}
\setlength{\parskip}{0.5em}

\setlist[itemize]{
    topsep=0.3em,
    itemsep=0.2em
}

\setlist[enumerate]{
    topsep=0.3em,
    itemsep=0.2em
}

\clubpenalty=10000
\widowpenalty=10000

\pagestyle{fancy}
\fancyhf{}
\fancyhead[L]{\nouppercase{\leftmark}}
\fancyhead[R]{\materia}
\fancyfoot[C]{\thepage}
\renewcommand{\headrulewidth}{0.4pt}
\renewcommand{\footrulewidth}{0pt}

\fancypagestyle{plain}{
    \fancyhf{}
    \fancyfoot[C]{\thepage}
    \renewcommand{\headrulewidth}{0pt}
}

% ============================================================
% COMANDOS MATEMÁTICOS
% ============================================================
\newcommand{\abs}[1]{\left|#1\right|}
\newcommand{\norm}[1]{\left\lVert#1\right\rVert}
\newcommand{\vect}[1]{\mathbf{#1}}
\newcommand{\deriv}[2]{\frac{d#1}{d#2}}
\newcommand{\pderiv}[2]{\frac{\partial#1}{\partial#2}}

% ============================================================
% ENTORNOS / CAJAS ACADÉMICAS
% ============================================================
\newtcolorbox{definition}[1]{
    enhanced,
    breakable,
    title={Definición: #1},
    fonttitle=\bfseries,
    colback=blue!3,
    colframe=blue!50!black,
    boxrule=0.8pt,
    arc=2mm
}

\newtcolorbox{important}{
    enhanced,
    breakable,
    title={Importante},
    fonttitle=\bfseries,
    colback=yellow!8,
    colframe=orange!70!black,
    boxrule=0.8pt,
    arc=2mm
}

\newtcolorbox{example}[1]{
    enhanced,
    breakable,
    title={Ejemplo: #1},
    fonttitle=\bfseries,
    colback=green!4,
    colframe=green!40!black,
    boxrule=0.8pt,
    arc=2mm
}

\newtcolorbox{formula}[1]{
    enhanced,
    breakable,
    title={Fórmula / Regla: #1},
    fonttitle=\bfseries,
    colback=purple!4,
    colframe=purple!50!black,
    boxrule=0.8pt,
    arc=2mm
}

\newtcolorbox{exercise}[1]{
    enhanced,
    breakable,
    title={Ejercicio: #1},
    fonttitle=\bfseries,
    colback=gray!5,
    colframe=gray!60!black,
    boxrule=0.8pt,
    arc=2mm
}

\newtcolorbox{note}{
    enhanced,
    breakable,
    title={Nota},
    fonttitle=\bfseries,
    colback=white,
    colframe=gray!50,
    boxrule=0.6pt,
    arc=2mm
}


% ============================================================
% DOCUMENT
% ============================================================
\begin{document}

% ------------------------------------------------------------
% PORTADA
% ------------------------------------------------------------
\begin{titlepage}
\thispagestyle{empty}

\begin{center}

\vspace*{1cm}

{\large \textsc{\universidad}}

\vspace{3cm}

{\Huge \textbf{\materia}}

\vspace{1cm}

{\LARGE \titulo}

\vspace{0.8cm}

\textit{Comprender a fondo cada institución jurídica es el primer paso para ejercerla con responsabilidad profesional.}

\vspace{1.2cm}

\rule{0.70\textwidth}{0.5pt}

\vspace{0.5cm}

{\large Apuntes de estudio}

\vspace{0.5cm}

\rule{0.70\textwidth}{0.5pt}

\vfill

{\large \autor}

\vspace{0.3cm}

{\large \anio}

\end{center}

\vfill

\begin{flushleft}
\textit{Este es un modesto aporte a los estudiantes de ``Derecho Notarial y Contratos'', no pretende sustituir el material oficial de la materia.}
\end{flushleft}

\end{titlepage}

% ------------------------------------------------------------
% ÍNDICE
% ------------------------------------------------------------
\tableofcontents

% ============================================================
% CAPÍTULO 1 — FUNDAMENTOS DEL FIDEICOMISO
% ============================================================
\chapter{Fundamentos del fideicomiso}

\section{Origen y naturaleza jurídica del fideicomiso}

El contrato de fideicomiso constituye una figura de enorme relevancia práctica en el ejercicio profesional del derecho y de la escribanía, aunque no configura un derecho real. Su tratamiento resulta prioritario dentro del programa de la materia por la frecuencia con la que se emplea en la práctica notarial y en el asesoramiento inmobiliario.

Esta figura ingresó al ordenamiento jurídico argentino con la sanción de la \textbf{Ley N.º 24.441}. Históricamente fue recibida con reservas por parte de la doctrina, al punto de que el propio Vélez Sarsfield no la incorporó al Código Civil originario. No obstante, la idea que subyace al fideicomiso es de raíz ancestral: consiste, esencialmente, en tener en cuenta al otro, a quien se le confía algo.

\begin{definition}{Fiducia}
Darle a alguien confiable algo. Cuando una persona entrega un bien en fiducia, lo hace porque confía en aquel a quien se lo entrega.
\end{definition}

El fideicomiso admite además una comparación con el derecho de superficie: mientras este último es un derecho real poco utilizado en la práctica, el fideicomiso, pese a encontrarse \textbf{muy en boga}, puede generar inconvenientes considerables si no se lo utiliza con el debido cuidado técnico.

\section{Riesgos y buen uso del contrato}

La desconfianza social hacia el fideicomiso proviene, en gran medida, de malas prácticas de mercado: existe la percepción difundida de que, al incorporar un bien a un fideicomiso, el aportante ya adquirió la propiedad correspondiente, sin advertir los riesgos que en realidad asume. El problema, por tanto, no radica en la figura en sí misma, sino en su utilización sin el debido asesoramiento técnico.

\begin{important}
El fideicomiso es un contrato que exige un análisis particular en cada caso: no debe copiarse un modelo genérico, sino redactarse a medida de la necesidad concreta del cliente. Aunque los sujetos que intervienen (fiduciante, fiduciario, beneficiario y fideicomisario) son siempre los mismos, la estructura contractual puede orientarse a finalidades muy distintas: de administración, de garantía, inmobiliarias o testamentarias, entre otras.
\end{important}

\begin{table}[H]\centering
\begin{tabularx}{\textwidth}{
    >{\raggedright\arraybackslash}p{0.30\textwidth}
    >{\raggedright\arraybackslash}X}
\toprule
\textbf{Preguntas / palabras clave} & \textbf{Notas / respuestas} \\
\midrule
\textbf{¿Qué es la fiducia?} &
Darle a alguien confiable algo: entregar un bien a quien se confía, con la expectativa de un manejo leal. \\
\textbf{¿Cuál es el origen normativo del fideicomiso en Argentina?} &
La Ley N.º 24.441, que introdujo la figura pese a las reservas históricas de la doctrina y a que Vélez Sarsfield no la había incorporado al Código Civil originario. \\
\textbf{¿Por qué genera desconfianza social el fideicomiso?} &
Por malas prácticas de mercado: se asume erróneamente que aportar un bien al fideicomiso equivale a adquirir la propiedad, sin advertir los riesgos que ello implica. \\
\textbf{¿Cuáles son los sujetos constantes del contrato de fideicomiso?} &
Fiduciante, fiduciario, beneficiario y fideicomisario, cuyas funciones se mantienen aunque la finalidad del contrato varíe (administración, garantía, inmobiliaria, testamentaria, entre otras). \\
\midrule
\multicolumn{2}{p{\dimexpr\textwidth-2\tabcolsep}}{
\textbf{Resumen:} El fideicomiso ingresó al derecho argentino mediante la Ley N.º 24.441 y se basa en la fiducia, es decir, en la confianza depositada en quien recibe un bien. No es un derecho real y su utilidad depende de que se redacte a medida de cada cliente, y no como un modelo genérico.} \\
\bottomrule
\end{tabularx}
\end{table}

% ============================================================
% CAPÍTULO 2 — SUJETOS Y PATRIMONIO FIDUCIARIO
% ============================================================
\chapter{Sujetos y patrimonio fiduciario}

\section{Fiduciante y fiduciario}

\begin{definition}{Fiduciante}
Es quien confía y aporta los bienes al fideicomiso: es el que confía, el que aporta.
\end{definition}

\begin{definition}{Fiduciario}
Es la persona que recibe los bienes transmitidos por el fiduciante. Constituye el punto más importante de esta unidad la naturaleza del derecho que el fiduciario adquiere sobre esos bienes.
\end{definition}

\begin{important}
El fiduciario no recibe los bienes como \textbf{dominio perfecto}, sino como \textbf{patrimonio fiduciario}. Esta distinción se vincula con los dominios imperfectos y constituye la idea central que debe recordarse sobre la posición del fiduciario: el patrimonio que recibe no le pertenece; lo recibe exclusivamente en su carácter fiduciario.
\end{important}

\begin{example}{El consorcio de propietarios}
En un informe de dominio referido a un departamento o a un terreno puede figurar como titular determinada persona. La respuesta técnicamente correcta no es ``soy la dueña'' sin más, sino ``lo recibí como dueña, pero como patrimonio fiduciario: no es mío''.

En la práctica, suele solicitarse el certificado de dominio a nombre del consorcio de propietarios de determinado inmueble. Sin embargo, el consorcio no es dueño del bien: el dominio se encuentra en cabeza del fiduciario, y quien vende lo hace en su carácter de fiduciario del consorcio, y no el consorcio mismo. La confusión entre estos planos dio lugar, en la práctica, a escrituras mal redactadas que debieron enmendarse porque el Registro las observaba, precisamente porque quien figuraba como vendedor no era el verdadero titular del dominio fiduciario.
\end{example}

\section{Beneficiario y fideicomisario}

\begin{definition}{Beneficiario}
Es quien recibe los beneficios del ejercicio de la administración fiduciaria. Puede ser el propio fiduciante o un tercero designado por él, pero no puede ser, a la vez, el fideicomisario.
\end{definition}

\begin{definition}{Fideicomisario}
Es quien recibe la propiedad de los bienes fideicomitidos una vez finalizado el contrato: el remanente. Ese remanente no necesariamente consiste en un inmueble, dado que este puede haberse distribuido ya al adjudicarse a los beneficiarios; puede tratarse, en cambio, de dinero remanente de la venta de las unidades, una gratificación, un IVA a favor no distribuido o un crédito no cobrado.
\end{definition}

\begin{important}
El fiduciario no puede ser, a la vez, el fideicomisario. Permitirlo generaría un conflicto de intereses estructural, dado que el fiduciario debe desempeñarse con lealtad, como \textbf{buen hombre de negocios} y como \textbf{buen padre de familia} —fórmula recogida por el propio Código—, sin tener un interés propio contrapuesto al de quienes confiaron en él.
\end{important}

\section{Protección del patrimonio fiduciario}

\begin{definition}{Separación e inmunidad patrimonial}
El patrimonio fiduciario no puede ser atacado por los acreedores del fiduciario, precisamente porque no le pertenece a título personal. Esa protección se extiende, en principio, a los acreedores del propio fiduciante: una vez que un bien ingresa al patrimonio fideicomitido, dicho patrimonio no puede ser agredido por los acreedores del fiduciante, salvo en el supuesto de nulidad del contrato.
\end{definition}

\begin{table}[H]\centering
\begin{tabularx}{\textwidth}{
    >{\raggedright\arraybackslash}p{0.30\textwidth}
    >{\raggedright\arraybackslash}X}
\toprule
\textbf{Preguntas / palabras clave} & \textbf{Notas / respuestas} \\
\midrule
\textbf{¿Qué recibe el fiduciario: dominio perfecto o patrimonio fiduciario?} &
Patrimonio fiduciario. No le pertenece a título personal; lo administra exclusivamente en función de su carácter fiduciario. \\
\textbf{Caso del consorcio: ¿quién vende el bien fideicomitido?} &
Quien vende es el fiduciario, en su carácter de tal, nunca el consorcio ni la sociedad, que no pueden ser titulares del dominio fiduciario. \\
\textbf{¿Quién puede ser beneficiario?} &
El propio fiduciante o un tercero designado por él; nunca, a la vez, el fideicomisario. \\
\textbf{¿Qué recibe el fideicomisario?} &
El remanente del patrimonio fideicomitido al finalizar el contrato: puede ser un bien, dinero, un crédito no cobrado o un IVA a favor. \\
\textbf{¿Por qué el fiduciario no puede ser fideicomisario?} &
Porque generaría un conflicto de intereses: el fiduciario debe actuar como buen hombre de negocios y buen padre de familia, sin interés propio contrapuesto al de quienes confiaron en él. \\
\textbf{¿Qué protección tiene el patrimonio fiduciario frente a los acreedores?} &
No puede ser atacado por los acreedores del fiduciario ni, en principio, por los del fiduciante, salvo que el contrato sea declarado nulo. \\
\midrule
\multicolumn{2}{p{\dimexpr\textwidth-2\tabcolsep}}{
\textbf{Resumen:} El fiduciario recibe los bienes como patrimonio fiduciario, no como dominio perfecto, y debe actuar con lealtad hacia fiduciante y beneficiarios. El fideicomisario recibe el remanente al finalizar el contrato y no puede coincidir con el fiduciario. La separación patrimonial protege esos bienes frente a los acreedores de ambas partes, salvo nulidad del contrato.} \\
\bottomrule
\end{tabularx}
\end{table}

% ============================================================
% CAPÍTULO 3 — TIPOS DE FIDEICOMISO SEGÚN SU FINALIDAD
% ============================================================
\chapter{Tipos de fideicomiso según su finalidad}

\section{Fideicomiso de acciones (sindicación)}

El paquete de control de una sociedad no equivale necesariamente a la titularidad del 50\,\% de su capital: en sociedades que cotizan en bolsa, quien reúne el 30\,\% puede, en la práctica, controlar la designación de directores o del consejo de vigilancia. A partir de este escenario resulta útil comparar dos alternativas para asegurar una decisión asamblearia común.

\begin{definition}{Pacto de sindicación sin fideicomiso}
Es válido entre las partes y les es oponible entre sí, pero no es oponible a la sociedad ni a terceros, salvo que la sociedad hubiera tomado conocimiento. Si un socio vota en la asamblea en sentido distinto al pactado, incumpliendo el pacto, la decisión asamblearia sigue siendo válida: solo generará sanciones por incumplimiento contractual entre las partes, sin afectar la validez del acto societario.
\end{definition}

\begin{definition}{Fideicomiso de acciones}
Los accionistas transmiten sus acciones a un fideicomiso; el fiduciario representa a esos fiduciantes-accionistas en la asamblea y vota en el sentido acordado previamente. De este modo, la decisión asamblearia queda directamente asegurada, porque quien vota es el fiduciario y no cada accionista individualmente.
\end{definition}

\begin{note}
% TODO: contenido incompleto o ambiguo en las notas originales
En relación con la reforma proyectada a la Ley General de Sociedades, el apunte original indica que el proyecto amplía la posibilidad de definir el objeto social por rubro de actividad, sin necesidad de detallar cada acto comprendido, en línea con la amplitud que ya presenta el objeto social de las SAS. Esa amplitud incide en la responsabilidad de los administradores por actos ``notoriamente extraños al objeto social'': si el objeto se define de manera amplia y todo lo realizado es lícito, prácticamente nada queda fuera de él, lo que reduce el margen para cuestionar la responsabilidad del directorio por actos ajenos a ese objeto. El apunte no especifica el estado o número de dicho proyecto de reforma.
\end{note}

\begin{important}
Como recomendación de ejercicio profesional aplicable a toda la práctica notarial: quien ejerce la abogacía o la escribanía no debe imponer su propia doctrina o pensamiento en los instrumentos que redacta. Debe considerar siempre la existencia de doctrina minoritaria, porque un tercero podría basarse en ella para atacar el acto en el futuro. El deber profesional consiste en cuidar al cliente y al tercero que confiará en el instrumento redactado, y no en hacer prevalecer la posición personal del profesional interviniente.
\end{important}

\section{Fideicomiso de administración}

Entre los bienes que pueden ser objeto de un fideicomiso se cuentan los inmuebles, los muebles, los derechos, la propiedad intelectual, las acciones y los créditos. El \textbf{fideicomiso de administración} es aquel en el cual el fiduciario administra un patrimonio ajeno con las facultades que el propio contrato le confiera; por ejemplo, administrar sin poder disponer, o administrar debiendo reunir previamente una asamblea de fiduciantes para ciertas decisiones, bajo pena de nulidad del acto.

\begin{example}{El dilema del fideicomiso de administración familiar}
Una persona transmite bienes a su hijo como fiduciario de administración, confiando en él, y prevé que a su fallecimiento el remanente pase a sus propios hijos (fideicomisarios), más allá de la legítima.

Este esquema plantea una tensión difícil de resolver: si el fiduciante desconfiara de sus hijos como fideicomisarios, debería preguntarse qué tan poca confianza tiene, entonces, en ese mismo hijo para desempeñarse como fiduciario. Como nadie sabe con certeza cuándo va a fallecer, no es posible prever con exactitud el momento en que el fiduciario debería transferir los bienes al fideicomisario, lo que obliga muchas veces a construir estructuras jurídicas cada vez más complejas y costosas, cuando la figura debería ser totalmente flexible.
\end{example}

\begin{note}
A modo de referencia comparada, el apunte original señala que en países como México y Estados Unidos el trust se encuentra integrado con naturalidad al patrimonio de prácticamente cualquier persona, y no solo de grandes patrimonios, principalmente por razones de practicidad y de costo, dado que en esos países los procesos sucesorios resultan sumamente costosos. En Argentina, el fideicomiso comenzó a funcionar recién a partir de la sanción de la Ley N.º 24.441.
\end{note}

\begin{important}
No todo fideicomiso es de ocultamiento. El riesgo se encuentra en quienes utilizan la figura de manera indebida, pero la figura en sí misma no es mala.
\end{important}

\section{Fideicomiso inmobiliario}

El fideicomiso inmobiliario admite una configuración sumamente amplia: el fiduciante puede aportar el terreno mismo, o bien dinero para que el fideicomiso, a través del fiduciario, adquiera el inmueble a fideicomitir. Cabe precisar un punto conceptual central: el fideicomiso no es sujeto de derecho; quien adquiere el inmueble es el fiduciario, en su carácter fiduciario, con el dinero aportado por los fiduciantes.

Según lo pactado en el contrato, el fiduciario podrá designar a quienes llevarán a cabo la obra, o bien los propios fiduciantes podrán haber contratado ya a esas personas, situación en la que puede producirse una superposición de roles.

\begin{example}{El comprador de la unidad ``a construir''}
Un tercero se entusiasma con un departamento en construcción dentro de un fideicomiso al costo y desea adquirirlo. Sin embargo, ese tercero no compra el departamento en sí: aporta dinero al fideicomiso para que la obra se construya, asumiendo el carácter de fiduciante.

Frente a la resistencia de algunos compradores a asumir ese rol, no existe alternativa jurídica distinta dentro de esa estructura: no puede instrumentarse un boleto de compraventa, porque el vínculo jurídico correcto es la \textbf{cesión de la posición contractual} que ya tenía el fiduciante original dentro del fideicomiso, con todos los derechos y también con todas las obligaciones (créditos y débitos) que esa posición conlleva.
\end{example}

\begin{example}{El caso de la cesión de posición contractual mal explicada}
Una persona que estaba comprando un departamento recibió un contrato de cesión de la posición contractual para su firma. Se advirtió que no debía firmarse ese documento tal como estaba redactado, porque, de existir una deuda dentro del fideicomiso, quien asumiera el carácter de fiduciante-beneficiaria por cesión pasaría a hacerse cargo de esa deuda. Ante la resistencia a modificar el texto, se sugirió incorporar una cláusula por la cual el cedente asumiera expresamente, al momento de la escritura, cualquier deuda existente en ese momento del fideicomiso.
\end{example}

\begin{important}
Cuando se cede la posición contractual dentro de un fideicomiso, se ceden tanto los créditos como los débitos existentes. Este extremo debe advertirse especialmente al momento de asesorar al cliente.
\end{important}

\begin{table}[H]\centering
\begin{tabularx}{\textwidth}{
    >{\raggedright\arraybackslash}p{0.30\textwidth}
    >{\raggedright\arraybackslash}X}
\toprule
\textbf{Preguntas / palabras clave} & \textbf{Notas / respuestas} \\
\midrule
\textbf{¿Qué ventaja tiene sindicar acciones mediante un fideicomiso, frente a un pacto simple?} &
El pacto simple es válido solo entre partes y no obliga a la sociedad; el fideicomiso de acciones asegura el voto porque es el fiduciario quien vota en la asamblea. \\
\textbf{Ética profesional: ¿debe el escribano imponer su propia doctrina?} &
No. Debe considerar siempre la doctrina minoritaria, porque un tercero podría basarse en ella para atacar el acto; su deber es cuidar al cliente y al tercero. \\
\textbf{¿Qué riesgo presenta el fideicomiso de administración familiar?} &
El conflicto entre la confianza depositada en el fiduciario y la eventual designación de fideicomisarios distintos, agravado por la imposibilidad de prever el momento del fallecimiento del fiduciante. \\
\textbf{Fideicomiso inmobiliario: ¿qué compra quien adquiere una unidad ``a construir''?} &
No compra el inmueble: aporta dinero al fideicomiso como fiduciante, o recibe una cesión de la posición contractual del fiduciante original. \\
\textbf{¿Qué implica la cesión de posición contractual en un fideicomiso inmobiliario?} &
Implica asumir tanto los créditos como los débitos (deudas) existentes en ese fideicomiso, no solo los derechos sobre la unidad. \\
\midrule
\multicolumn{2}{p{\dimexpr\textwidth-2\tabcolsep}}{
\textbf{Resumen:} Según su finalidad, el fideicomiso puede emplearse para sindicar acciones y asegurar el voto en asambleas, para administrar patrimonios familiares o para financiar y llevar adelante obras inmobiliarias. En todos los casos, el fiduciario actúa a nombre propio pero en interés ajeno, y la cesión de posiciones contractuales dentro del fideicomiso conlleva tanto derechos como obligaciones.} \\
\bottomrule
\end{tabularx}
\end{table}

% ============================================================
% CAPÍTULO 4 — FIDEICOMISO, SUCESIÓN Y LEGÍTIMA
% ============================================================
\chapter{Fideicomiso, sucesión y legítima}

\section{El fideicomiso frente a la sucesión}

La cuestión de si el patrimonio fideicomitido ingresa o no a la sucesión del fiduciante no es unánime en la doctrina.

\begin{important}
Si el fideicomiso está bien redactado, el patrimonio fideicomitido no entra en la sucesión, porque se trata de un patrimonio que ya fue transmitido. No se trata de una postura unánime: existen abogados que no comparten este criterio porque, de aceptarlo, quedan sin la sucesión como vía de reclamo; otros sostienen la misma posición aquí expuesta.
\end{important}

Si el fiduciante designó fideicomisarios (por ejemplo, sus propios hijos), el bien no debería reingresar a la sucesión al finalizar el fideicomiso, porque ya cuenta con un destinatario contractualmente previsto. En cambio, si el contrato no previó un fideicomisario y el bien está destinado a ``volver'' al propio fiduciante fallecido, el fideicomiso se encuentra mal redactado y, en ese supuesto, el bien remanente sí ingresa a la sucesión.

\begin{definition}{Plazo del fideicomiso y sucesores universales}
El fideicomiso tiene un plazo máximo de treinta años. Para prever el fallecimiento del fiduciante dentro de ese plazo, es posible designar como fideicomisario no solo a una persona determinada, sino también a ``quienes resulten sus sucesores universales'', evitando así que, ante la falta de previsión, el bien termine reingresando a la sucesión.
\end{definition}

\section{Legítima, colación y porción disponible}

Mientras exista la legítima tal como está regulada actualmente, todo aquello que la afecte resulta colacionable: si el fiduciante designa como fideicomisario a uno solo de sus hijos, en detrimento de los demás herederos forzosos, esos otros herederos podrán reclamar la colación de lo que exceda la porción disponible del causante.

\begin{important}
Todo aquello que afecte la legítima es colacionable mientras se mantenga la regulación actual de esta institución.
\end{important}

\begin{definition}{Cálculo de la porción disponible}
La porción disponible no debe fijarse en un testamento como un porcentaje rígido (por ejemplo, ``el 33\,\%''), sino como ``la porción disponible al momento en que se produzca el fallecimiento'', ya que ese porcentaje puede variar con el tiempo según eventuales reformas legales. Si al momento del fallecimiento la porción disponible aumentó, por ejemplo, al 50\,\%, la disposición testamentaria alcanzará ese 50\,\%.
\end{definition}

\begin{table}[H]\centering
\begin{tabularx}{\textwidth}{
    >{\raggedright\arraybackslash}p{0.30\textwidth}
    >{\raggedright\arraybackslash}X}
\toprule
\textbf{Preguntas / palabras clave} & \textbf{Notas / respuestas} \\
\midrule
\textbf{¿El patrimonio fideicomitido ingresa a la sucesión del fiduciante?} &
No, si el fideicomiso está bien redactado y tiene fideicomisario designado; si no lo tiene, el bien remanente sí reingresa a la sucesión. \\
\textbf{¿Cuál es el plazo máximo del fideicomiso?} &
Treinta años; por eso conviene prever como fideicomisario no solo a una persona, sino también a ``quienes resulten sus sucesores universales''. \\
\textbf{¿Qué es colacionable en materia de legítima?} &
Todo aquello que afecte la legítima de los herederos forzosos, como designar fideicomisario a un solo hijo en perjuicio de sus hermanos. \\
\textbf{¿Cómo debe fijarse la porción disponible en un testamento?} &
Como la porción disponible vigente al momento del fallecimiento, no como un porcentaje fijo, porque la ley puede modificar ese porcentaje con el tiempo. \\
\midrule
\multicolumn{2}{p{\dimexpr\textwidth-2\tabcolsep}}{
\textbf{Resumen:} Un fideicomiso correctamente redactado, con fideicomisario designado, evita que el patrimonio fideicomitido reingrese a la sucesión del fiduciante, aunque este criterio no es unánime en la doctrina. La legítima, sin embargo, condiciona esa planificación: todo lo que la afecte es colacionable, y la porción disponible debe calcularse según su vigencia al momento del fallecimiento, no como un porcentaje fijo.} \\
\bottomrule
\end{tabularx}
\end{table}

% ============================================================
% CAPÍTULO 5 — CASOS DE LA PRÁCTICA PROFESIONAL
% ============================================================
\chapter{Casos de la práctica profesional}

\section{Testamento, secreto profesional y conflicto familiar}

La voluntad testamentaria puede desatar conflictos familiares imprevistos, por lo que corresponde actuar con especial cuidado en la redacción y el otorgamiento de testamentos.

\begin{example}{El testamento y la disputa familiar}
Una persona, plenamente capaz al momento de testar y casada en segundas nupcias, otorgó testamento dejando constancia expresa de que, si bien tenía hijos —dos residentes en el exterior y una en el país—, quien se ocupaba permanentemente de su cuidado era su esposa. El testador manifestó su voluntad de que sus hijos comprendieran esa decisión al leer el testamento.

Con el paso del tiempo, los hijos regresaron al país y advirtieron que su padre presentaba un cuadro de deterioro. Le reclamaron entonces al escribano que les informara qué había firmado su padre. Por el secreto profesional, el escribano no podía revelar el contenido de la escritura, y sugirió a los hijos que fuera el propio padre quien decidiera mostrarles o no ese documento.

El padre, tiempo después, contactó nuevamente al escribano, ya en presencia de los hijos, y este le preguntó directamente si recordaba haber firmado un testamento. El padre respondió que sí, pero que en realidad no había sido su iniciativa: fue su esposa quien lo llevó a la escribanía en varias oportunidades para ese fin.

Finalmente, el padre —queriendo evitar el conflicto familiar— buscó junto al escribano una salida intermedia: dado que no podía revocar unilateralmente lo ya dispuesto sin generar una situación conflictiva con su esposa, se acordó redactar un compromiso entre las partes sin valor legal frente a terceros, pero con valor moral entre ellas, por el cual la esposa se comprometía a no reclamar, al momento del fallecimiento del marido, la porción disponible que le había sido asignada. Ese compromiso se instrumentaría luego como un acta de manifestación —que sí constituye un instrumento público—, aunque sin hacer referencia a la renuncia en sí, ``porque no tiene ningún valor'' jurídico vinculante como tal.
\end{example}

\begin{important}
Como recomendación práctica dirigida al ejercicio de la función notarial: ante situaciones de otorgamiento testamentario, corresponde hacer comparecer previamente a los interesados, y no instrumentar el testamento por vía telefónica, de modo que quede claro qué es lo efectivamente querido por el testador.
\end{important}

\begin{note}
El secreto profesional del escribano impone un límite claro frente a la familia del testador: no puede revelar el contenido de una escritura, aun cuando los herederos forzosos lo requieran, y solo puede sugerir que sea el propio otorgante quien decida compartir esa información.
\end{note}

\section{Incapacidad sobreviniente y riesgos de la administración}

La misma inestabilidad de la voluntad humana frente a las presiones familiares puede presentarse también en el ámbito de la administración de un fideicomiso, retomando lo ya analizado sobre el fideicomiso de administración.

\begin{example}{El fiduciante que necesita revisar la administración}
Un fiduciante transmitió todos sus bienes a un fiduciario de su confianza para que los administrara, contando con una renta mensual para su propio sustento. Con el tiempo, ese fiduciante considera que la renta asignada ya no le alcanza para vivir de la mejor manera posible, mientras que, desde la perspectiva del fiduciario, la administración es objetivamente correcta: las cuentas están invertidas en fondos de inversión o en plazo fijo, y el fiduciario no retiene más que los honorarios pactados en el contrato de fideicomiso.

A este conflicto de percepciones se suma una segunda variable: la posibilidad de que el propio fiduciante pierda, con el tiempo, la capacidad de evaluar razonablemente esa administración. A modo de ejemplo, se mencionó el caso de una persona que, tras un episodio de presión arterial alta, fue trasladada a un centro de cuidado alejado del entorno social de la persona, de modo que resultaba muy difícil para sus allegados visitarla con frecuencia.
\end{example}

\begin{important}
Al redactar un fideicomiso de administración debe preverse expresamente qué ocurre cuando el fiduciante, al no encontrarse ya en condiciones psíquicas para evaluar la administración, necesita revisar o modificar el contrato. Existen numerosas situaciones posibles que, con frecuencia, no se encuentran previstas.
\end{important}

\begin{table}[H]\centering
\begin{tabularx}{\textwidth}{
    >{\raggedright\arraybackslash}p{0.30\textwidth}
    >{\raggedright\arraybackslash}X}
\toprule
\textbf{Preguntas / palabras clave} & \textbf{Notas / respuestas} \\
\midrule
\textbf{Caso del testamento: ¿qué límite impone el secreto profesional?} &
El escribano no puede revelar a los hijos el contenido de la escritura; solo puede sugerir que sea el propio testador quien decida mostrarla. \\
\textbf{¿Qué valor tiene un ``compromiso entre partes'' que renuncia informalmente a la porción disponible?} &
No tiene valor legal vinculante frente a terceros; puede instrumentarse como acta de manifestación (instrumento público), pero sin mencionar la renuncia en sí. \\
\textbf{¿Qué recomendación práctica se formula para el otorgamiento de testamentos?} &
Hacer comparecer previamente a los interesados y no instrumentar el testamento por vía telefónica, para que quede claro qué es lo efectivamente querido por el testador. \\
\textbf{¿Qué riesgo adicional presenta la administración fiduciaria frente a la incapacidad sobreviniente?} &
Que el fiduciante pierda, con el tiempo, la capacidad de evaluar razonablemente la administración, sin que el contrato haya previsto un mecanismo de revisión para ese supuesto. \\
\midrule
\multicolumn{2}{p{\dimexpr\textwidth-2\tabcolsep}}{
\textbf{Resumen:} Los casos reales de la práctica profesional muestran que los mayores riesgos del fideicomiso y del testamento no provienen de las figuras en sí, sino de una redacción imprevisora: la confusión entre titular y fiduciario, la cesión mal explicada de una posición contractual, el conflicto familiar desatado por un testamento y el deterioro de la capacidad de un fiduciante exigen previsión contractual expresa y resguardo del secreto profesional.} \\
\bottomrule
\end{tabularx}
\end{table}

% ============================================================
% CAPÍTULO 6 — CIERRE Y TRABAJO PARA LA PRÓXIMA CLASE
% ============================================================
\chapter{Cierre y trabajo para la próxima clase}

\section{Fideicomiso en garantía y testamentario: consigna}

Quedan pendientes de desarrollo dos tipos de fideicomiso considerados de particular importancia: el \textbf{fideicomiso en garantía} y el \textbf{fideicomiso testamentario}, que se desarrollarán en un encuentro posterior.

\begin{exercise}{Consigna para el próximo encuentro}
A partir de lo ya visto sobre el fideicomiso de administración y el fideicomiso inmobiliario, deberán elaborarse ejemplos propios centrados no en los usos posibles de cada figura, sino en los \textbf{conflictos concretos} que pueden generarse en cada una de ellas, de modo que, al momento de redactar el contrato, dichos conflictos no pasen inadvertidos.

A título orientador: en un fideicomiso inmobiliario donde no se previó adecuadamente la figura del fideicomisario, podría verse afectada la legítima; o bien, quien recibió una cesión de posición contractual podría luego pretender desconocer una deuda anterior —por ejemplo, honorarios de un constructor o de una medianera— alegando que ``solo compró'' el departamento y no la posición contractual completa.
\end{exercise}

\begin{important}
Comprender los conflictos posibles de cada tipo de fideicomiso —y no solo sus usos— permite proteger adecuadamente al cliente y decidir con fundamento qué incluir y qué no incluir en el contrato.
\end{important}

% ============================================================
% CUADRO CORNELL GENERAL DE REPASO
% ============================================================
\chapter{Cuadro Cornell general de repaso}

\begin{longtable}{
    >{\raggedright\arraybackslash}p{0.30\textwidth}
    >{\raggedright\arraybackslash}p{0.60\textwidth}}
\toprule
\textbf{Preguntas / palabras clave} & \textbf{Notas / respuestas} \\
\midrule
\endfirsthead
\toprule
\textbf{Preguntas / palabras clave} & \textbf{Notas / respuestas} \\
\midrule
\endhead

\textbf{¿Qué es la fiducia?} &
Darle a alguien confiable algo: entregar un bien a quien se confía, con la expectativa de un manejo leal. \\
\textbf{¿Cuál es el origen normativo del fideicomiso en Argentina?} &
La Ley N.º 24.441, que introdujo la figura pese a las reservas históricas de la doctrina y a que Vélez Sarsfield no la había incorporado. \\
\textbf{¿Qué recibe el fiduciario: dominio perfecto o patrimonio fiduciario?} &
Patrimonio fiduciario. No le pertenece a título personal; lo administra en función de su carácter fiduciario. \\
\textbf{Caso del consorcio: ¿quién vende el bien fideicomitido?} &
Quien vende es el fiduciario, en su carácter de tal, nunca el consorcio ni la sociedad, que no pueden ser titulares del dominio fiduciario. \\
\textbf{¿Quién puede ser beneficiario?} &
El propio fiduciante o un tercero designado por él; nunca, a la vez, el fideicomisario. \\
\textbf{¿Qué recibe el fideicomisario?} &
El remanente del patrimonio fideicomitido al finalizar el contrato: puede ser un bien, dinero, un crédito no cobrado o un IVA a favor. \\
\textbf{¿Por qué el fiduciario no puede ser fideicomisario?} &
Porque generaría un conflicto de intereses: el fiduciario debe actuar como buen hombre de negocios, sin interés propio contrapuesto al de quienes confiaron en él. \\
\textbf{¿Qué protección tiene el patrimonio fiduciario frente a acreedores?} &
No puede ser atacado por los acreedores del fiduciario ni, en principio, por los del fiduciante, salvo que el contrato sea declarado nulo. \\
\textbf{¿Qué ventaja tiene sindicar acciones mediante un fideicomiso, frente a un pacto simple?} &
El pacto simple es válido solo entre partes y no obliga a la sociedad; el fideicomiso de acciones asegura el voto porque es el fiduciario quien vota en la asamblea. \\
\textbf{Ética profesional: ¿debe el escribano imponer su propia doctrina?} &
No. Debe considerar siempre la doctrina minoritaria, porque un tercero podría basarse en ella para atacar el acto; su deber es cuidar al cliente y al tercero. \\
\textbf{¿Qué riesgo presenta el fideicomiso de administración familiar?} &
El conflicto entre la confianza depositada en el fiduciario y la eventual designación de fideicomisarios distintos, agravado por la imposibilidad de prever el momento del fallecimiento del fiduciante. \\
\textbf{Fideicomiso inmobiliario: ¿qué compra quien adquiere una unidad ``a construir''?} &
No compra el inmueble: aporta dinero al fideicomiso como fiduciante, o recibe una cesión de la posición contractual del fiduciante original. \\
\textbf{¿Qué implica la cesión de posición contractual en un fideicomiso inmobiliario?} &
Implica asumir tanto los créditos como los débitos (deudas) existentes en ese fideicomiso, no solo los derechos sobre la unidad. \\
\textbf{¿El patrimonio fideicomitido ingresa a la sucesión del fiduciante?} &
No, si el fideicomiso está bien redactado y tiene fideicomisario designado; si no lo tiene, el bien remanente sí reingresa a la sucesión. \\
\textbf{¿Cuál es el plazo máximo del fideicomiso?} &
Treinta años; por eso conviene prever como fideicomisario no solo a una persona, sino también a ``quienes resulten sus sucesores universales''. \\
\textbf{¿Qué es colacionable en materia de legítima?} &
Todo aquello que afecte la legítima de los herederos forzosos, como designar fideicomisario a un solo hijo en perjuicio de sus hermanos. \\
\textbf{¿Cómo debe fijarse la porción disponible en un testamento?} &
Como la porción disponible vigente al momento del fallecimiento, no como un porcentaje fijo, porque la ley puede modificar ese porcentaje con el tiempo. \\
\textbf{Caso del testamento: ¿qué límite impone el secreto profesional?} &
El escribano no puede revelar a los hijos el contenido de la escritura; solo puede sugerir que sea el propio testador quien decida mostrarla. \\
\textbf{¿Qué valor tiene un ``compromiso entre partes'' que renuncia informalmente a la porción disponible?} &
No tiene valor legal vinculante frente a terceros; puede instrumentarse como acta de manifestación (instrumento público), pero sin mencionar la renuncia en sí. \\
\textbf{Consigna final: ¿qué debía preparar el curso para la clase siguiente?} &
Ejemplos propios de fideicomiso de administración e inmobiliario centrados en los conflictos que pueden generar, no en sus usos posibles. \\
\bottomrule
\end{longtable}

\begin{important}
\textbf{Síntesis final.} El fideicomiso es una figura contractual —no un derecho real— construida sobre la confianza que el fiduciante deposita en el fiduciario, quien recibe los bienes exclusivamente como patrimonio fiduciario, separado tanto de su propio dominio como de los acreedores de ambas partes. Esa separación patrimonial es, precisamente, lo que permite emplear el fideicomiso con finalidades muy distintas: asegurar el voto en una asamblea societaria mediante la sindicación de acciones, administrar un patrimonio familiar, financiar y llevar adelante una obra inmobiliaria, o planificar el destino de los bienes tras el fallecimiento del fiduciante, sin que se trate nunca de un contrato ``modelo'' para copiar, sino de una estructura que debe redactarse a medida de cada cliente. Los casos reales analizados —la confusión entre el consorcio y su fiduciario, la cesión mal explicada de una posición contractual, el testamento que reabrió un conflicto familiar, el deterioro de capacidad de un fiduciante— muestran que los mayores riesgos de esta figura no provienen del instituto en sí, sino de una redacción imprevisora. De allí la advertencia central de la materia: quien ejerce la abogacía o la escribanía debe anticipar los conflictos posibles —no solo los usos posibles— de cada fideicomiso, cuidando siempre al cliente y a los terceros que, en el futuro, confiarán en el instrumento redactado.
\end{important}

\end{document}
